\documentclass[12pt]{article}
\usepackage[T1]{fontenc}
\usepackage[utf8]{inputenc}
\usepackage{lmodern}
\usepackage{textcomp}
\usepackage{enumitem}
\usepackage{geometry}
\geometry{margin=1in}
\begin{document}
\begin{center}
Answer Key - Health Sciences - Graduate Level Assessment 6 \
\small (Answers are ordered exactly as in the question paper.)
\end{center}

\section*{Multiple Choice}
\begin{enumerate}[label=\arabic*.]
\item \textbf{Answer:} C
\\ \textit{Options:} A) The end-product of glycolysis in all cells is glucose; B) Glycolysis doesn't require any enzymes; C) In glycolysis, glucose is cleaved into two three-carbon products; D) Glucose metabolism only occurs in the presence of oxygen; E) In glycolysis, glucose is converted into two molecules of fructose
\\ \textit{Note:} The correct answer is accurate because glycolysis is the metabolic pathway that converts glucose into two molecules of pyruvate, each containing three carbon atoms, through a series of enzymatic reactions.
\item \textbf{Answer:} E
\\ \textit{Options:} A) Vitamin C; B) Vitamin E; C) Vitamin B$_{12}$; D) Thiamin; E) Riboflavin
\\ \textit{Note:} Riboflavin, as a precursor for the coenzyme FAD (flavin adenine dinucleotide), is essential for the fatty acid β-oxidation process, where it plays a crucial role in the oxidation of fatty acids to produce energy.
\item \textbf{Answer:} A
\\ \textit{Options:} A) Personality; B) Skin elasticity; C) Physical strength; D) Male hormone levels; E) Hair color
\\ \textit{Note:} Personality is considered a discontinuous change over the adult years because it can undergo significant transformations influenced by life experiences, relationships, and personal development, rather than evolving gradually.
\item \textbf{Answer:} B
\\ \textit{Options:} A) Poor muscle development; B) Increased risk of mortality; C) Goitre; D) Increased risk of obesity; E) Increased risk of heart disease
\\ \textit{Note:} Vitamin A deficiency in children is linked to an increased risk of mortality due to its critical role in maintaining immune function, vision, and overall health, making them more susceptible to infections and other health complications.
\item \textbf{Answer:} B
\\ \textit{Options:} A) Sunscreen; B) Vaccines; C) Antivirals; D) Pain Relievers; E) Hand Sanitizer
\\ \textit{Note:} Vaccines are designed to stimulate the immune system and provide immunity against specific diseases, making them one of the most effective methods for disease prevention.
\end{enumerate}

\section*{True / False}
\begin{enumerate}[label=\arabic*.]
\setcounter{enumi}{5}
\item \textbf{Answer:} False
\\ \textit{Options:} A) True; B) False
\\ \textit{Note:} Familial hypercholesterolaemia is caused by mutations in the low-density lipoprotein receptor (LDLR) gene, not in the gene that encodes apolipoprotein E.
\item \textbf{Answer:} True
\\ \textit{Options:} A) True; B) False
\\ \textit{Note:} Phosphocreatine is predominantly found in the cytoplasm of muscle cells because it serves as a rapid source of energy by donating a phosphate group to adenosine diphosphate (ADP) to regenerate adenosine triphosphate (ATP) during high-intensity muscular activity.
\item \textbf{Answer:} True
\\ \textit{Options:} A) True; B) False
\\ \textit{Note:} A stoma should ideally be sited in the rectus sheath abdominus because this location provides better support, minimizes complications, and facilitates easier access for care and management.
\item \textbf{Answer:} False
\\ \textit{Options:} A) True; B) False
\\ \textit{Note:} Astroviruses primarily infect the intestinal tract and are associated with gastroenteritis, rather than causing significant replication in the liver or leading to hepatic infection.
\item \textbf{Answer:} True
\\ \textit{Options:} A) True; B) False
\\ \textit{Note:} Hepatitis C was first identified in 1989 through the use of molecular biology techniques, specifically the detection of the virus's RNA, rather than through clinical observation.
\end{enumerate}

\section*{Long Form Response}
\begin{enumerate}[label=\arabic*.]
\setcounter{enumi}{10}
\item \textbf{Answer:} Food balance sheets, while useful for tracking the supply and distribution of food within a country, have significant limitations in providing a comprehensive understanding of food security and nutrition. Firstly, they often focus solely on quantitative measures, neglecting qualitative aspects such as the nutritional value and safety of food, which are critical for assessing dietary adequacy. Moreover, food balance sheets typically aggregate data at national levels, obscuring regional disparities and the varying accessibility of food among different population groups. Additionally, they do not account for non-food factors influencing food security, such as socio-economic conditions, cultural preferences, and the impact of food waste. For instance, a country may exhibit a balanced food supply in terms of calories but still face malnutrition issues due to poor dietary diversity or inequitable food distribution. Thus, while food balance sheets provide a foundational overview, they fall short of delivering a nuanced understanding of the complexities surrounding food security and nutrition.
\\ \textit{Note:} The answer correctly identifies that food balance sheets primarily provide quantitative data on food supply and distribution, failing to account for qualitative factors like nutritional value and regional disparities, which are essential for a comprehensive understanding of food security and nutrition.
\item \textbf{Answer:} Various treatments for depression in older adults, including pharmacotherapy, psychotherapy, and lifestyle interventions, have demonstrated varying degrees of effectiveness. Antidepressant medications, particularly selective serotonin reuptake inhibitors (SSRIs), are commonly prescribed and can significantly alleviate symptoms of depression. Psychotherapeutic approaches, such as cognitive-behavioral therapy (CBT), also show promise by addressing negative thought patterns and enhancing coping strategies. In contrast, mental stimulation, while beneficial for cognitive health, may not be considered an effective intervention for depression specifically, as it often fails to address the underlying emotional and psychological factors contributing to depressive symptoms. This limitation suggests that while mental stimulation can improve cognitive function, it may not provide the targeted emotional support needed to effectively treat depression in older adults.
\\ \textit{Note:} correct because it identifies and evaluates the main treatment options for depression in older adults, highlighting the effectiveness of pharmacotherapy and psychotherapy while setting the stage for a critical analysis of the limitations of mental stimulation as an intervention.
\end{enumerate}

\vfill
\noindent\textit{End of Answer Key}
\end{document}